\section{Abelian varieties of Weil type}\label{sec:weiltype}

We fix the notion of Weil type, construct the Weil line, prove that it
consists of Hodge classes, and record the two facts about it that the
obstruction theorems use: its dimension, and its position relative to the
subring generated by divisor classes. The standard references are \cite{BL04}
for abelian varieties and \cite{Voi02} for Hodge theory. Apart from the
Riemann relations, the Lefschetz theorem on $(1,1)$ classes and the
description of the N\'eron--Severi group by symmetric endomorphisms, every
result of the section is proved here or in \Cref{sec:characters}.

\subsection{Weil type}

Let $K=\QQ(\sqrt{-d})$ be an imaginary quadratic field, with $d$ a squarefree
positive integer, and let $\bbar{\phantom{x}}$ denote the nontrivial element
of $\Gal(K/\QQ)$.

\begin{definition}[Weil type]\label{def:weiltype}
Let $A$ be a complex abelian variety, let
$\iota\colon K\hookrightarrow\End^{0}(A)$ be an embedding, and let $\eta$ be a
polarisation class on $A$. The triple $(A,\iota,\eta)$ is of
\emph{$(K,-1,n)$-Weil type} if
\begin{enumerate}[label=\textup{(\roman*)},nosep]
\item $\dim A=2n$;
\item the endomorphism $\iota(\sqrt{-d})$ acts on the tangent space
  $H^{0}(A,\Omega^{1}_{A})^{\vee}\cong H^{1,0}(A)^{\vee}$ with the two
  eigenvalues $\pm\sqrt{-d}$ each occurring with multiplicity $n$;
\item the Rosati involution attached to $\eta$ restricts on
  $\iota(K)\subset\End^{0}(A)$ to complex conjugation, that is
  $\iota(\tau)^{\dagger}=\iota(\bbar{\tau})$ for all $\tau\in K$.
\end{enumerate}
\end{definition}

Condition (ii) is the balanced signature condition. It separates Weil type
from the unbalanced case, in which the eigenvalues $\pm\sqrt{-d}$ occur with
different multiplicities, and it is the condition under which the Weil line
consists of Hodge classes (\Cref{rem:balancecount}). Condition (iii) says that
the polarisation is compatible with the $K$-action, and the middle entry $-1$
of $(K,-1,n)$ records the sign in this compatibility: the Rosati involution
restricts to conjugation on $\iota(K)$, not to the identity. By
\Cref{lem:rosatiauto} the sign is forced once $\iota(K)$ is stable under the
Rosati involution. Some authors suppress the sign and speak of $K$-Weil type
of dimension $2n$.

\begin{lemma}[$H^{1}$ is a $K$-space]\label{lem:Kspace}
Let $(A,\iota,\eta)$ be of $(K,-1,n)$-Weil type and put $V=H^{1}(A,\QQ)$. Then
$\iota$ makes $V$ a $K$-vector space of dimension $2n$, and
$\dim_{\QQ}V=4n$.
\end{lemma}

\begin{proof}
The algebra $\End^{0}(A)$ acts $\QQ$-linearly on $V$ by pullback. Pullback
reverses the order of composition, but $K$ is commutative, so
$\tau\mapsto\iota(\tau)^{*}$ is a unital ring homomorphism
$K\to\End_{\QQ}(V)$. Since $K$ is a field, $V$ becomes a $K$-vector space,
and since $\dim_{\QQ}V=2\dim A=4n$, we get $\dim_{K}V=4n/[K:\QQ]=2n$.
\end{proof}

\begin{notation}\label{not:eigen}
Write $V_{\CC}=V\otimes_{\QQ}\CC$. Since $\iota(\sqrt{-d})^{2}=-d$ on $V$,
the operator $\iota(\sqrt{-d})\otimes1$ on $V_{\CC}$ has the two
eigenvalues $\pm\sqrt{-d}$, and
\begin{equation}\label{eq:Vsplit}
  V_{\CC}=V_{+}\oplus V_{-},
  \qquad
  V_{\pm}=\ker\bigl(\iota(\sqrt{-d})\mp\sqrt{-d}\bigr),
\end{equation}
with $\dim_{\CC}V_{\pm}=2n$ each, because $V$ is free of rank $2n$ over $K$ and
$K\otimes_{\QQ}\CC\cong\CC\times\CC$. Complex conjugation on $V_{\CC}$, coming
from the real structure $V\otimes\RR$, interchanges $V_{+}$ and $V_{-}$. For
general $\tau=a+b\sqrt{-d}\in K$ the operator $\iota(\tau)$ acts on $V_{+}$ by
the scalar $\tau$ and on $V_{-}$ by $\bbar{\tau}$.
\end{notation}

\begin{remark}\label{rem:tautbar}
The two scalars in \Cref{not:eigen} are different functions of $\tau$. For
$\tau\notin\QQ$ the operator $\iota(\tau)$ is $\QQ$-linear on $V$ but not a
scalar; it becomes a scalar on each piece only after extension to $\CC$ and
splitting. This is the source of the characters studied in
\Cref{sec:characters}. A class defined over $\QQ$ can be an eigenvector of
every $\iota(\tau)$ only with an eigenvalue in $\QQ$; the function
$\Nm(\tau)=\tau\bbar{\tau}$ is $\QQ$-valued, and $\tau^{2n}$ is not
(\Cref{rem:rightcharacter}).
\end{remark}

\subsection{The Hodge structure}

Write the Hodge decomposition as $V_{\CC}=V^{1,0}\oplus V^{0,1}$, with
$V^{1,0}=H^{1,0}(A)$ and $V^{0,1}=H^{0,1}(A)$. The $K$-action preserves it,
because $\iota(\tau)$ is induced by an endomorphism of $A$ and is therefore a
morphism of Hodge structures. The eigenspaces of an operator that preserves
both summands are the direct sums of its eigenspaces on the two summands, so
each of $V_{+}$ and $V_{-}$ splits: with $V^{p,q}_{\pm}=V_{\pm}\cap V^{p,q}$,
\begin{equation}\label{eq:bisplit}
  V_{\pm}=V_{\pm}^{1,0}\oplus V_{\pm}^{0,1}.
\end{equation}

\begin{lemma}[Balanced signature]\label{lem:balanced}
For $(A,\iota,\eta)$ of $(K,-1,n)$-Weil type,
\[
  \dim_{\CC}V_{+}^{1,0}=\dim_{\CC}V_{-}^{1,0}=n,
  \qquad
  \dim_{\CC}V_{+}^{0,1}=\dim_{\CC}V_{-}^{0,1}=n .
\]
\end{lemma}

\begin{proof}
By \Cref{def:weiltype}(ii), $\iota(\sqrt{-d})$ acts on the tangent space with
each eigenvalue of multiplicity $n$. The tangent space is dual to
$H^{1,0}(A)=V^{1,0}$, and the action of $\iota(\sqrt{-d})$ on $V^{1,0}$ by
pullback is the transpose of its action on the tangent space, so it has the
same eigenvalues with the same multiplicities; that is,
$\dim V^{1,0}_{+}=\dim V^{1,0}_{-}=n$. Since $\dim_{\CC}V_{\pm}=2n$ by
\Cref{not:eigen} and $V_{\pm}=V^{1,0}_{\pm}\oplus V^{0,1}_{\pm}$, the
remaining two dimensions are $n$ each.
\end{proof}

\subsection{The Weil line}

\begin{definition}[Weil line]\label{def:weilline}
Let $(A,\iota,\eta)$ be of $(K,-1,n)$-Weil type and $V=H^{1}(A,\QQ)$. The
\emph{Weil line} is
\[
  \HW(A)\;=\;\Bigl(\textstyle\bigwedge^{2n}_{\CC}V_{+}\oplus
  \bigwedge^{2n}_{\CC}V_{-}\Bigr)\cap H^{2n}(A,\QQ),
\]
the intersection being taken inside
$H^{2n}(A,\CC)=\bigwedge^{2n}_{\CC}V_{\CC}$.
\end{definition}

\begin{proposition}[Structure of the Weil line]\label{prop:weilline}
Let $(A,\iota,\eta)$ be of $(K,-1,n)$-Weil type. Then
\begin{enumerate}[label=\textup{(\roman*)},nosep]
\item $\HW(A)$ is a $\QQ$-subspace of $H^{2n}(A,\QQ)$ of dimension $2$, stable
  under $\iota(K)$, and free of rank one as a $K$-module;
\item $\HW(A)\subset H^{n,n}(A)$, so every element of $\HW(A)$ is a Hodge class;
\item $\HW(A)$ is canonically isomorphic to $\bigwedge^{2n}_{K}V$ as a
  $K$-module.
\end{enumerate}
\end{proposition}

\begin{proof}
(i) Write $W_{\CC}=\bigwedge^{2n}V_{+}\oplus\bigwedge^{2n}V_{-}$. Each summand
is one-dimensional over $\CC$ because $\dim_{\CC}V_{\pm}=2n$, so
$\dim_{\CC}W_{\CC}=2$. The splitting \eqref{eq:Vsplit} is the eigenspace
decomposition of the operator $\iota(\sqrt{-d})$, which is defined over $\QQ$,
so every automorphism of $\CC$, acting on $V_{\CC}$ through the second
factor, either preserves $V_{+}$ and $V_{-}$ or interchanges them, according
to its restriction to $K$; complex conjugation interchanges them. Hence
$W_{\CC}$ is stable under $\mathrm{Aut}(\CC)$, and by Galois descent it is
defined over $\QQ$: the space $\HW(A)=W_{\CC}\cap H^{2n}(A,\QQ)$ has
$\dim_{\QQ}\HW(A)=2$ and $\HW(A)\otimes\CC=W_{\CC}$. It is stable under $\iota(K)$ because
$\iota(\tau)$ preserves each of $V_{\pm}$, hence each summand. A
two-dimensional $\QQ$-space with a faithful $K$-action is free of rank one
over $K$.

(ii) By \eqref{eq:bisplit} and \Cref{lem:balanced}, the line
$\bigwedge^{2n}V_{+}$ equals
$\bigwedge^{n}V^{1,0}_{+}\otimes\bigwedge^{n}V^{0,1}_{+}$, which is of Hodge
bidegree $(n,n)$. The same computation applies to $\bigwedge^{2n}V_{-}$.
Hence $W_{\CC}\subset H^{n,n}(A)$ and so
$\HW(A)\subset H^{2n}(A,\QQ)\cap H^{n,n}(A)=\Hdg^{n}(A)$.

(iii) Since $K\otimes_{\QQ}\CC\cong\CC\times\CC$, with the two factors
corresponding to the two embeddings of $K$, we have
$V\otimes_{\QQ}\CC=V\otimes_{K}(K\otimes_{\QQ}\CC)=V_{+}\oplus V_{-}$, and
exterior powers over $K$ commute with this base change:
\[
  \Bigl(\textstyle\bigwedge^{2n}_{K}V\Bigr)\otimes_{\QQ}\CC
  \;\cong\;\textstyle\bigwedge^{2n}_{\CC}V_{+}\oplus\bigwedge^{2n}_{\CC}V_{-}
  \;=\;W_{\CC}\;\subset\;H^{2n}(A,\CC).
\]
Every map in this chain is compatible with the action of $\mathrm{Aut}(\CC)$
through the second factor, which on the right permutes $V_{+}$ and $V_{-}$ as in (i).
By Galois descent the composite is induced by an injective $\QQ$-linear map
$\bigwedge^{2n}_{K}V\to H^{2n}(A,\QQ)$, which commutes with $\iota(K)$ because
each map in the chain does. Its image is a $\QQ$-subspace of
$W_{\CC}\cap H^{2n}(A,\QQ)=\HW(A)$ of dimension two, hence all of $\HW(A)$.
\end{proof}

\begin{remark}\label{rem:onedim}
Some authors call a generator of $\HW(A)$ \emph{the} Weil class and speak of a
one-dimensional space; the difference is the usual one between a $K$-line and
its underlying $\QQ$-plane. Nothing below depends on the convention, since
the obstruction theorems concern the whole of $\HW(A)$.
\end{remark}

\subsection{The Weil line and the divisor subring}

The Weil line is not accounted for by divisor classes. We record the
statement in the form in which it is used and derive it from the character
computation of \Cref{sec:characters}.

\begin{definition}\label{def:divisorsubring}
The \emph{divisor subring} $D^{*}(A)\subset H^{*}(A,\QQ)$ is the
$\QQ$-subalgebra generated by $\NS(A)\otimes\QQ\subset H^{2}(A,\QQ)$.
\end{definition}

Every element of $D^{*}(A)$ is algebraic, by the Lefschetz theorem on $(1,1)$
classes and the fact that the algebraic classes form a subring. The Hodge
conjecture for $A$ in degree $n$ is therefore equivalent to the algebraicity
of a complement of $D^{2n}(A)$ in $\Hdg^{n}(A)$. For a very general member of
a Weil family the Weil line is such a complement, by \Cref{thm:hdgcount}; at
special members the divisor subring can grow and meet the Weil line
(\Cref{rem:special}).

\begin{proposition}\label{prop:notpoly}
Let $(A,\iota,\eta)$ be of $(K,-1,n)$-Weil type with
$\End^{0}(A)=\iota(K)$, and suppose $n\ge1$. Then
\[
  \HW(A)\cap D^{2n}(A)=0 .
\]
\end{proposition}

\begin{proof}
By \Cref{prop:weilline}(i) and \Cref{thm:character}, the space
$\HW(A)\otimes\CC$ is the sum of the $\Kx$-eigenspaces for the characters
$\tau\mapsto\tau^{2n}$ and $\tau\mapsto\bbar{\tau}^{2n}$. The space
$\NS(A)\otimes\QQ$ is isomorphic to the space of elements of $\End^{0}(A)$
fixed by the Rosati involution \cite[Section~5.2]{BL04}. Since
$\End^{0}(A)=\iota(K)$ and the Rosati involution restricts to conjugation on
$\iota(K)$, the fixed elements are those of $\iota(\QQ)$, so
$\NS(A)\otimes\QQ=\QQ\,\eta$. By \Cref{lem:NSchar} the class $\eta$ is a
$\Kx$-eigenclass for the norm character $\Nm$, so every element of
$D^{2n}(A)=\QQ\,\eta^{n}$ is an eigenclass for $\Nm^{n}$. By
\Cref{lem:charsdiffer} the characters $\tau\mapsto\tau^{2n}$ and
$\tau\mapsto\bbar{\tau}^{2n}$ both differ from $\tau\mapsto\Nm(\tau)^{n}$.
Eigenvectors for pairwise distinct characters are linearly independent, so
$\HW(A)\otimes\CC$ meets the $\Nm^{n}$-eigenspace, which contains
$D^{2n}(A)$, only in zero.
\end{proof}

The mechanism behind the proof is that of \Cref{rem:rightcharacter}: $\Kx$
can act on a rational class through a character only if the character is
$\QQ$-valued, and the $\QQ$-valued characters occurring in the cohomology
are the powers of the norm. The Weil line has rank one over $K$, not over $\QQ$, and as a
$K$-line of exterior degree $2n$ it carries the $2n$-th power of $\tau$
itself. \Cref{thm:character} and \Cref{lem:NSchar,lem:charsdiffer} are proved
in \Cref{sec:characters}, and nothing there uses \Cref{prop:notpoly}, so the
forward reference is not circular.
